\chapter{Datasets and Data Engineering}
\label{ch:datasets}

\section{CICIDS2017}

CICIDS2017~\cite{Sharafaldin2018,CIC2017} is a five-day capture from the
Canadian Institute for Cybersecurity testbed, with approximately
2.83~million flow records labelled across 15 classes. After cleaning
(Section~\ref{sec:cleaning}) we retain 2{,}313{,}150 flows with 73
features each. Attack classes are concentrated by day:

\begin{table}[H]
    \centering
    \caption{CICIDS2017 attack class distribution by day.}
    \label{tab:cicids-by-day}
    \small
    \begin{tabular}{lll}
        \toprule
        \textbf{Day} & \textbf{Role (day-split)} & \textbf{Attack types present} \\
        \midrule
        Monday & Train & Benign only \\
        Tuesday & Train & Benign, FTP-Patator, SSH-Patator \\
        Wednesday & Train & Benign, DoS Hulk, GoldenEye, slowloris, Slowhttptest, Heartbleed \\
        Thursday & Validation & Benign, Web-BF, Web-XSS, Web-SQLi, Infiltration \\
        Friday & Test & Benign, DDoS, PortScan, Bot \\
        \bottomrule
    \end{tabular}
\end{table}

Friday's test classes (DDoS, PortScan, Bot) never appear in
Monday-to-Wednesday training. This is a structural reason that
multi-class macro-F1 collapses on the day split: a supervised classifier
cannot output a label it has never seen.

\begin{figure}[H]
    \centering
    \includegraphics[width=0.95\textwidth]{day_attack_distribution.png}
    \caption{CICIDS2017 attack flows per day, coloured by attack type
    (log y-axis). Shaded regions mark the day-based train/val/test
    split.}
    \label{fig:day-attack-distribution}
\end{figure}

The dataset also exhibits extreme class imbalance: 85\% of flows are
benign, the most common attack class (DoS Hulk) has 172{,}688 flows, and
the rarest (Heartbleed) has only 11 --- a 220{,}000:1 ratio. We mitigate
imbalance with class-weighting strategies documented in
Chapter~\ref{ch:method} and use macro F1 as the headline metric.

\begin{figure}[H]
    \centering
    \includegraphics[width=0.8\textwidth]{class_imbalance.png}
    \caption{CICIDS2017 class distribution on a log scale.}
    \label{fig:class-imbalance}
\end{figure}

\section{UNSW-NB15 and the Need for Two Datasets}

UNSW-NB15~\cite{Moustafa2015,UNSW2015} is a 2.54-million-row capture from
the IXIA PerfectStorm platform with nine attack categories plus Normal.
We use the official train/test split with 175{,}341 test rows. UNSW
provides a useful cross-dataset comparison because it was generated by a
different team on a different testbed with a substantially different
attack catalogue, and the binary label distribution is more balanced
(attack rate $\approx$ 31\% versus CICIDS' 14\%).

A single-dataset claim is not really a claim about machine learning --- it
is a claim about that one dataset. By benchmarking on two independent
corpora we can distinguish robust patterns (which generalise) from
dataset artefacts (which do not). The rank-reversal between the two
datasets reported in Chapter~\ref{ch:crossdataset} is the main empirical
justification for the ``no universal winner'' conclusion.

\section{Data Cleaning Pipeline}
\label{sec:cleaning}

Raw CSVs from CIC contain duplicate rows, infinity values from divisions
by zero, and inconsistent column whitespace. The
\texttt{src/prepare\_data.py} module performs the following
deterministically: strip whitespace from column names; replace
$\pm\infty$ with NaN and drop NaN rows; deduplicate by hashing the full
feature-plus-label vector; add day-of-week and source-file columns; and
persist to a Parquet file of approximately 200~MB.

A key step is \textit{leakage filtering}. Engelen et al.~\cite{Engelen2021}
showed that several CICIDS2017 papers achieved their headline numbers
because the four rate features --- Flow Bytes/s, Flow Packets/s, Forward
and Backward Packets/s --- leak the label through their denominators. We
drop these four columns explicitly and the \texttt{src/leakage\_check.py}
module enforces their absence. Removing them alone lowers stratified
random-forest macro F1 by approximately four percentage points.
\label{sec:leakage}

\section{Train, Validation, and Test Splits}

Three split policies, each progressively more demanding. The
\textbf{stratified} split is a row-wise stratified partition preserving
class proportions across train, validation, and test partitions --- the
policy used in most prior work, producing the optimistic numbers. The
\textbf{day-based} split corresponds directly to the day-of-capture
column: train on Monday--Wednesday, validate on Thursday, test on Friday.
This is the realistic protocol for a deployed IDS. The
\textbf{leave-one-day-out (LODO)} split holds out one day for testing,
uses the remaining days for training and threshold-tuning validation, and
rotates through all five days. LODO is the strongest temporal
generalisation test because it exercises the model on five disjoint
held-out distributions.

A stratified split lets the model see Friday traffic during training,
memorise Friday-specific patterns, and report them as generalisation. A
day-based split forbids this, and the 0.4-macro-F1 gap between the two
policies is the memorisation tax. This pattern matches Pendlebury et
al.~\cite{Pendlebury2019}, who showed for malware classification that
random splits over-state deployment performance and proposed time-aware
evaluation as the default protocol --- a recommendation we adopt here.
